\documentclass[10pt,a4paper]{article}
% BTS CIEL 2e année — Overleaf : pdfLaTeX (deux compilations)
% Mise en page : chaque TP tient sur une seule page. Le corrigé est dans
% tp_routage_dynamique_corrige.tex.
\usepackage[utf8]{inputenc}
\usepackage[T1]{fontenc}
\usepackage[french,provide=*]{babel}
\usepackage{lmodern,microtype}
\usepackage[a4paper,hmargin=1.4cm,top=1.3cm,bottom=1.4cm,includehead,headheight=13pt,headsep=4mm,footskip=6mm]{geometry}
\usepackage{xcolor,tabularx,booktabs,array,enumitem,fancyhdr,lastpage}
\usepackage{tikz}
\usetikzlibrary{arrows.meta,positioning,calc}
\usepackage[most]{tcolorbox}
\usepackage{listings,hyperref,titlesec}
\definecolor{vert}{HTML}{1C322D}
\definecolor{creme}{HTML}{F8F3EE}
\definecolor{or}{HTML}{EBB552}
\definecolor{rose}{HTML}{F1CDBE}
\definecolor{rouge}{HTML}{B03A2E}
\definecolor{gris}{HTML}{697672}
\hypersetup{colorlinks=true,linkcolor=vert,urlcolor=vert}
\setlength{\parindent}{0pt}
\setlength{\parskip}{2pt}
\setlist[enumerate]{leftmargin=*,itemsep=4pt}
\setlist[itemize]{leftmargin=*,itemsep=2pt}
\renewcommand{\arraystretch}{1.15}
\pagestyle{fancy}\fancyhf{}
\fancyhead[L]{\small\color{vert}BTS CIEL — Routage}
\fancyhead[R]{\small TP en cinq séances}
\fancyfoot[L]{\small Ahmed Bachir Diedhiou}
\fancyfoot[R]{\small\thepage/\pageref{LastPage}}
\renewcommand{\headrulewidth}{0.4pt}
\titleformat{\section}{\Large\bfseries\color{vert}}{\thesection.}{.6em}{}
\titleformat{\subsection}{\large\bfseries\color{vert}}{\thesubsection}{.55em}{}
\titlespacing*{\section}{0pt}{0pt}{4pt}
\tcbset{compact/.style={enhanced,breakable,arc=2mm,boxsep=1mm,left=2mm,right=2mm,top=.5mm,bottom=.5mm,toptitle=.3mm,bottomtitle=.3mm,before skip=2.5mm,after skip=4pt}}
% Le titre de chaque encadré ouvre son texte, pour gagner une ligne.
\newtcolorbox{probleme}{compact,colback=creme,colframe=or,before upper={{\bfseries\color{vert}Situation-problème —}\enspace}}
\newtcolorbox{notion}{compact,colback=creme,colframe=vert,before upper={{\bfseries\color{vert}Comprendre avant d'agir —}\enspace}}
\newtcolorbox{bilan}{compact,colback=white,colframe=vert,before upper={{\bfseries\color{vert}Synthèse à compléter —}\enspace}}
% Lignes de réponse : l'espace d'écriture est au-dessus de chaque trait.
\newcommand{\lignes}[1][1]{\par\foreach \k in {1,...,#1}{\vspace{3.5mm}\noindent\rule{\linewidth}{.3pt}\par}\vspace{-1mm}}
\newcommand{\question}[2][1]{\par\textbf{\color{vert}Question —} #2\lignes[#1]}
\newcommand{\etape}[1]{\par\vspace{1pt}{\bfseries\color{vert}#1\enspace}}
\newcommand{\duree}[1]{\hfill{\normalsize\mdseries #1}}
\lstset{basicstyle=\ttfamily\footnotesize,backgroundcolor=\color{creme},frame=single,rulecolor=\color{gris},columns=fullflexible,keepspaces=true,breaklines=true,showstringspaces=false,xleftmargin=2mm,xrightmargin=2mm,aboveskip=3pt,belowskip=3pt}

% Un même schéma sert aux trois états du projet. Les étiquettes montrent
% volontairement les IP exactes des interfaces, les ports, les PC et les LAN.
% #1 : état (1, 2 ou 3) ; #2 : largeur du schéma.
\newcommand{\topologie}[2]{%
\par{\centering\resizebox{#2}{!}{%
\begin{tikzpicture}[font=\sffamily\scriptsize,>=Latex,yscale=.8,
 routeur/.style={circle,draw=vert,fill=vert,text=white,minimum size=11mm,font=\bfseries},
 pc/.style={draw=vert,fill=creme,rounded corners=1mm,align=center,minimum width=29mm,minimum height=10mm},
 commut/.style={draw=vert,fill=rose,rounded corners=1mm,align=center,minimum width=18mm,minimum height=7mm}]
\node[routeur] (r1) at (0,0) {R1};
\node[routeur] (r2) at (8,0) {R2};
\node[commut] (s1) at (0,-1.65) {SW1};
\node[commut] (s2) at (8,-1.65) {SW2};
\node[pc] (p1) at (0,-3.05) {PC Paris\\192.168.1.10/24\\GW 192.168.1.1};
\node[pc] (p2) at (8,-3.05) {PC Lyon\\192.168.2.10/24\\GW 192.168.2.1};
\draw[thick] (r1)--node[right,align=left]{G0/0\\192.168.1.1/24}(s1);
\draw[thick] (s1)--(p1);
\draw[thick] (r2)--node[right,align=left]{G0/0\\192.168.2.1/24}(s2);
\draw[thick] (s2)--(p2);
\draw[line width=1.5pt,or] (r1)--node[above]{10.0.12.0/30} node[pos=.12,below,align=center]{G0/1\\10.0.12.1} node[pos=.88,below,align=center]{G0/1\\10.0.12.2}(r2);
\ifnum#1>1
\node[routeur] (r3) at (16,0) {R3};
\node[commut] (s3) at (16,-1.65) {SW3};
\node[pc] (p3) at (16,-3.05) {PC Marseille\\192.168.3.10/24\\GW 192.168.3.1};
\draw[thick] (r3)--node[right,align=left]{G0/0\\192.168.3.1/24}(s3);
\draw[thick] (s3)--(p3);
\draw[line width=1.5pt,or] (r2)--node[above]{10.0.23.0/30} node[pos=.12,below,align=center]{G0/2\\10.0.23.1} node[pos=.88,below,align=center]{G0/1\\10.0.23.2}(r3);
\fi
\ifnum#1>2
\draw[line width=1.6pt,or] (r1.north) .. controls (2.5,2.1) and (13.5,2.1) .. (r3.north)
 node[pos=.5,above]{Liaison de secours~: 10.0.13.0/30};
\node[anchor=south east,align=right] at (r1.north west) {R1 G0/2\\10.0.13.1};
\node[anchor=south west,align=left] at (r3.north east) {R3 G0/2\\10.0.13.2};
\fi
\end{tikzpicture}}\par}\smallskip}

\begin{document}
\begin{titlepage}\pagecolor{vert}\color{creme}
\vspace*{3cm}{\Large BTS CIEL — Deuxième année}\par\vspace{1cm}
{\Huge\bfseries Du routage statique\\[3mm]au routage dynamique}\par\vspace{8mm}
{\Large Un TP Packet Tracer en cinq séances, de RIP à OSPF}\par\vspace{1cm}
{\color{or}\rule{\linewidth}{2pt}}\par\vspace{1cm}
{\large 9 heures — situations-problèmes, manipulation, observation, explication}\par
\vfill
\begin{tcolorbox}[colback=creme,colframe=or]
\color{vert}\textbf{Question directrice~:} comment permettre aux routeurs d'apprendre les chemins et de s'adapter aux changements du réseau~?
\end{tcolorbox}
\vfill Ahmed Bachir Diedhiou
\end{titlepage}\nopagecolor\color{black}
\tableofcontents\clearpage

\section{Le réseau HexaTech~: projet commun}
HexaTech possède d'abord deux sites (Paris et Lyon), puis ouvre Marseille. Vous configurerez la même infrastructure tout au long des cinq séances.
\begin{center}\begin{tabularx}{\linewidth}{lXr}\toprule
Séance 1 & Deux sites et routes statiques & 1 h 30\\
Séance 2 & Troisième site et limites du statique & 1 h 30\\
Séance 3 & Découverte et configuration de RIPv2 & 2 h\\
Séance 4 & Redondance et convergence & 1 h 30\\
Séance 5 & Découverte et configuration d'OSPF & 2 h 30\\\bottomrule\end{tabularx}\end{center}
\textbf{Matériel~:} Packet Tracer, trois routeurs Cisco 2911, trois switches 2960, trois PC, câbles Ethernet. Chaque routeur utilise G0/0 pour le LAN~; les autres interfaces sont précisées sur les schémas.\par
\textbf{Règle de travail~:} enregistrer une copie distincte à chaque séance. Les tests se font d'abord en mode Realtime, puis en mode Simulation lorsque demandé.\par
\textbf{Commandes IOS~:} sauf mention contraire, les commandes de configuration se saisissent après \texttt{enable} puis \texttt{configure terminal}~; les commandes \texttt{show} se saisissent en mode privilégié (\texttt{Router\#}).

% ---------------------------------------------------------------- Séance 1
\clearpage\section[Séance 1 — Deux sites doivent communiquer]{Séance 1 — Deux sites doivent communiquer\duree{1 h 30}}
\begin{probleme}Le bureau de Paris doit échanger des données avec le bureau de Lyon. Chaque LAN fonctionne, mais les deux PC n'arrivent pas à communiquer. Pourquoi~?\end{probleme}
\topologie{1}{.56\linewidth}
\textbf{Objectifs~:} configurer les interfaces, tester la connectivité, lire la table de routage, créer les routes statiques nécessaires.

\etape{Étape 1 — Construire et adresser.}
Placez deux routeurs 2911, deux switches 2960 et deux PC, reliez-les et adressez les PC comme sur le schéma. Sur R1 comme sur R2, configurez \textbf{G0/0} (LAN) et \textbf{G0/1} (liaison entre routeurs), puis activez-les~:
\begin{lstlisting}
enable
configure terminal
interface gigabitEthernet 0/0
 ip address ADRESSE MASQUE
 no shutdown
exit
\end{lstlisting}
\question{Pourquoi les interfaces G0/0 et G0/1 d'un même routeur n'ont-elles pas des adresses dans le même réseau~?}

\etape{Étape 2 — Vérifier les interfaces.}
Sur les deux routeurs, lancez \texttt{show ip interface brief}.
\question{Quels états doivent apparaître dans les colonnes \texttt{Status} et \texttt{Protocol} pour une liaison opérationnelle~?}

\etape{Étape 3 — Constater le problème.}
Depuis PC Paris, exécutez \texttt{ping 192.168.1.1}, puis \texttt{ping 192.168.2.10}.
\question[2]{Quel test fonctionne~? Formulez une hypothèse expliquant l'échec de l'autre test.}

\etape{Étape 4 — Vérifier l'hypothèse.}
Sur R1 puis R2, exécutez \texttt{show ip route}. Repérez les réseaux précédés de \texttt{C} (connecté) et de \texttt{L} (adresse locale d'interface).
\question{R1 possède-t-il une route vers \texttt{192.168.2.0/24}~? Et R2 vers \texttt{192.168.1.0/24}~?}

\etape{Étape 5 — Corriger avec des routes statiques.}
Sur R1, complétez puis saisissez \texttt{ip route RESEAU\_DESTINATION MASQUE IP\_PROCHAIN\_SAUT}~; faites l'équivalent sur R2.
\question{Pourquoi faut-il prévoir une route dans chaque sens~?}
Refaites le ping entre les PC et consultez \texttt{show ip route}.
\question{Quelle lettre distingue maintenant les nouvelles routes dans la table~?}
\begin{bilan}Une route statique est \dotfill\\Un routeur connaît initialement \dotfill\\Le prochain saut désigne \dotfill\end{bilan}

% ---------------------------------------------------------------- Séance 2
\clearpage\section[Séance 2 — L'entreprise ouvre Marseille]{Séance 2 — L'entreprise ouvre Marseille\duree{1 h 30}}
\begin{probleme}Un troisième site ouvre à Marseille. Le responsable annonce ensuite vingt nouvelles implantations. La méthode utilisée précédemment pourra-t-elle être conservée sans difficulté~?\end{probleme}
\topologie{2}{\linewidth}

\etape{Étape 1 — Ajouter le troisième site.}
Reprenez une copie du fichier de la séance 1. Ajoutez R3, SW3 et PC Marseille. Sur R2, configurez \textbf{G0/2} vers R3~; sur R3, configurez \textbf{G0/1} vers R2 et \textbf{G0/0} vers son LAN. Attribuez toutes les IP figurant sur le schéma.
\question{Quels nouveaux réseaux apparaissent dans la table de R2~?}

\etape{Étape 2 — Tester avant de configurer les routes.}
Depuis PC Paris, lancez \texttt{ping 192.168.3.10}. Consultez ensuite les tables des trois routeurs.
\question[2]{Quels routeurs ne connaissent pas encore de chemin vers un LAN distant~? Lesquels leur manquent~?}

\etape{Étape 3 — Compléter les routes statiques.}
Ajoutez les routes nécessaires sur R1, R2 et R3, en tenant compte des routes de la séance 1. Testez Paris–Marseille, Lyon–Marseille et les trajets retour.
\question[2]{Combien de nouvelles routes avez-vous créées~? Quelles erreurs de configuration seraient possibles~?}

\etape{Étape 4 — Prendre du recul.}
L'entreprise prévoit désormais vingt sites supplémentaires.
\question[2]{Quelles tâches l'administrateur devra-t-il répéter à chaque extension~?}
\question[2]{Imaginez un moyen pour que les routeurs apprennent eux-mêmes les réseaux connus de leurs voisins.}
\begin{bilan}La principale limite du routage statique dans un grand réseau est \dotfill\\Nous aurions besoin que les routeurs \dotfill\end{bilan}

% ---------------------------------------------------------------- Séance 3
\clearpage\section[Séance 3 — Découvrir et configurer RIPv2]{Séance 3 — Découvrir et configurer RIPv2\duree{2 h}}
\begin{probleme}Le service informatique veut réduire la maintenance manuelle des routes. Comment R1 pourrait-il apprendre le réseau de Marseille grâce aux annonces de R2~?\end{probleme}
\topologie{2}{.9\linewidth}
\begin{notion}Un routeur peut annoncer à ses voisins les réseaux qu'il sait atteindre~; les voisins utilisent ces annonces pour mettre à jour leur table~: c'est le \textbf{routage dynamique}. \textbf{RIP} (Routing Information Protocol) est un protocole à \emph{vecteur de distance}~; RIPv2 annonce aussi les masques. Sa métrique est le \textbf{nombre de sauts}~: une métrique de 16 signifie « inaccessible ».\end{notion}

\etape{Étape 1 — Prédire la métrique.}
R3 possède le LAN Marseille~: R2 l'apprend en métrique 1, R1 en métrique 2.
\question{Pourquoi R1 doit-il envoyer les paquets à R2 pour atteindre Marseille dans cette topologie~?}

\etape{Étape 2 — Préparer une expérience correcte.}
Dupliquez le fichier de la séance 2. Sur chaque routeur, consultez \texttt{show ip route static}, puis supprimez les routes statiques des séances 1 et 2 avec \texttt{no ip route RESEAU MASQUE IP\_PROCHAIN\_SAUT}.
\question{Pourquoi faut-il retirer ces routes avant d'étudier les routes apprises par RIP~?}

\etape{Étape 3 — Configurer RIP sur R1.}
Saisissez, dans l'ordre, les commandes du tableau~:\par\smallskip
{\small\begin{tabularx}{\linewidth}{>{\ttfamily}p{.37\linewidth}X}\toprule
\normalfont\textbf{Commande}&\textbf{Rôle}\\\midrule
router rip&Active le processus RIP.\\
version 2&Utilise RIPv2, dont les annonces incluent les masques.\\
network 192.168.1.0&Active RIP sur l'interface du LAN Paris.\\
network 10.0.0.0&Active RIP sur les interfaces du réseau majeur 10.0.0.0/8.\\
no auto-summary&Désactive le résumé automatique aux frontières des réseaux majeurs.\\
passive-interface gigabitEthernet 0/0&Pas d'annonces RIP vers les PC~; le LAN reste annoncé.\\\bottomrule
\end{tabularx}}\par
\textbf{Attention~:} \texttt{network} ne crée aucune route~: il désigne les interfaces RIP (réseau majeur, d'où \texttt{10.0.0.0}).
\question{Pourquoi le LAN Paris est-il annoncé, alors que l'interface G0/0 est passive~?}

\etape{Étape 4 — Reproduire en comprenant.}
Sur R2 puis R3, configurez RIP en adaptant le réseau du LAN.
\question{Quelles commandes sont identiques sur les trois routeurs~? Lesquelles changent et pourquoi~?}

\etape{Étape 5 — Vérifier l'apprentissage.}
Sur R1, exécutez \texttt{show ip protocols} et \texttt{show ip route}, puis testez le ping Paris–Marseille dans les deux sens. Extrait attendu~:
\begin{lstlisting}
R    192.168.3.0/24 [120/2] via 10.0.12.2, 00:00:12, GigabitEthernet0/1
\end{lstlisting}
\question{Interprétez la lettre R, les valeurs 120 et 2, puis l'adresse 10.0.12.2.}

\etape{Étape 6 — Voir RIP fonctionner.}
En \textbf{Simulation}, filtrez \textbf{RIP} et ouvrez une annonce (toutes les 30 s, UDP 520, vers 224.0.0.9).
\question{Quelles informations sont transmises~? À quoi servent-elles à un routeur voisin~?}
\begin{bilan}RIP~: famille \makebox[5cm]{\dotfill}~; métrique \dotfill\\Lettre d'une route RIP dans la table~: \makebox[4cm]{\dotfill}\end{bilan}

% ---------------------------------------------------------------- Séance 4
\clearpage\section[Séance 4 — Une liaison tombe en panne]{Séance 4 — Une liaison tombe en panne\duree{1 h 30}}
\begin{probleme}Une liaison directe Paris–Marseille est créée pour disposer de plusieurs chemins. Que se passe-t-il si elle tombe en panne~?\end{probleme}
\topologie{3}{\linewidth}

\etape{Étape 1 — Ajouter le chemin direct.}
Reprenez le fichier de la séance 3. Reliez \textbf{R1 G0/2} à \textbf{R3 G0/2}, configurez les IP indiquées sur le schéma, puis \texttt{no shutdown}. Les interfaces en 10.0.13.0/30 sont déjà couvertes par \texttt{network 10.0.0.0} dans RIP.
\question{Combien de chemins physiques relient maintenant Paris à Marseille~?}

\etape{Étape 2 — Observer le trajet normal.}
Attendez la convergence puis, depuis PC Paris, exécutez \texttt{tracert 192.168.3.10}.
\question[2]{Quel chemin est utilisé~? Expliquez ce choix à partir de la métrique RIP.}

\etape{Étape 3 — Provoquer et observer une panne.}
Sur R1, désactivez G0/2~:
\begin{lstlisting}
configure terminal
interface gigabitEthernet 0/2
 shutdown
\end{lstlisting}
Après convergence, consultez \texttt{show ip route}, puis relancez le ping et le traceroute.
\question[2]{Par quel routeur les paquets passent-ils maintenant~? La communication est-elle nécessairement ininterrompue pendant la convergence~?}
\begin{notion}La \textbf{convergence} est le processus par lequel les routeurs prennent connaissance d'un changement et mettent à jour leurs routes. Le routage dynamique ne crée pas de nouvelle liaison~: pour contourner une panne, un chemin alternatif doit déjà exister.\end{notion}

\etape{Étape 4 — Rétablir la liaison.}
Sur R1 G0/2~: \texttt{no shutdown}. Attendez la convergence et observez la table de R1.
\question{Pourquoi le chemin préféré change-t-il à nouveau~?}
\begin{bilan}La redondance signifie \dotfill\\La convergence correspond à \dotfill\end{bilan}

% ---------------------------------------------------------------- Séance 5
\clearpage\section[Séance 5 — Découvrir et configurer OSPF]{Séance 5 — Découvrir et configurer OSPF\duree{2 h 30}}
\begin{probleme}Le chemin qui passe par le moins de routeurs n'est pas toujours le plus intéressant. Comment choisir selon le coût des liaisons plutôt que selon le seul nombre de sauts~?\end{probleme}
\topologie{3}{.72\linewidth}
\begin{notion}\textbf{OSPF} (\emph{Open Shortest Path First}) est un protocole à \textbf{état de liens}~: les routeurs découvrent leurs voisins (messages \textbf{Hello}), échangent la topologie de leur aire, puis calculent les chemins de coût minimal (algorithme SPF de Dijkstra). Ici, une seule aire~: \textbf{area 0}.\end{notion}

\etape{Étape 1 — Faire une hypothèse.}
Coûts prévus~: R1–R2 = 10~; R2–R3 = 10~; R1–R3 = 30. Le coût du LAN de R3, commun aux deux chemins, ne change pas le choix.
\question{Comparez les coûts des trajets R1–R3 et R1–R2–R3~: quel chemin OSPF devrait-il préférer vers Marseille~?}

\etape{Étape 2 — Configurer OSPF sur R1.}
Rétablissez toutes les interfaces, supprimez RIP partout (\texttt{no router rip}), puis saisissez sur R1~:\par\smallskip
{\small\begin{tabularx}{\linewidth}{>{\ttfamily}p{.37\linewidth}X}\toprule
\normalfont\textbf{Commande}&\textbf{Rôle}\\\midrule
router ospf 1&Active le processus OSPF 1 (numéro local, propre au routeur).\\
network 192.168.1.0 0.0.0.255 area 0&Place l'interface du LAN de R1 dans l'aire 0.\\
network 10.0.12.0 0.0.0.3 area 0&Place l'interface de R1 vers R2 dans l'aire 0.\\
network 10.0.13.0 0.0.0.3 area 0&Place l'interface de R1 vers R3 dans l'aire 0.\\
passive-interface gigabitEthernet 0/0&Pas de voisinage OSPF sur le LAN, qui reste annoncé.\\\bottomrule
\end{tabularx}}\par
\textbf{Masque générique (wildcard)} = \texttt{255.255.255.255 - masque} (/24~: \texttt{0.0.0.255}~; /30~: \texttt{0.0.0.3}).

\etape{Étape 3 — Adapter aux autres routeurs.}
Sur R2 et R3, adaptez les trois lignes \texttt{network} aux réseaux de leurs interfaces~; G0/0 reste passive.
\question{Quelles interfaces doivent établir des voisinages OSPF~? Pourquoi pas G0/0~?}

\etape{Étape 4 — Observer les voisins et les routes.}
Sur R1, exécutez \texttt{show ip ospf neighbor}, puis \texttt{show ip route}.
\question{Combien de voisins R1 découvre-t-il~? Quelle lettre désigne les routes OSPF~? Recopiez la route vers \texttt{192.168.3.0/24}.}

\etape{Étape 5 — Appliquer les coûts.}
En mode interface, sur R1~: \texttt{ip ospf cost 10} sur G0/1, \texttt{ip ospf cost 30} sur G0/2~; R2~: 10 sur G0/1 et G0/2~; R3~: 10 sur G0/1, 30 sur G0/2. Vérifiez avec \texttt{show ip ospf interface}, relisez la route vers \texttt{192.168.3.0/24} et lancez \texttt{tracert 192.168.3.10} depuis PC Paris.
\question[2]{Qu'est-ce qui a changé dans la route~? Que représentent les deux nombres entre crochets~? Le chemin choisi confirme-t-il votre hypothèse~?}

\etape{Étape 6 — Couper une liaison.}
Coupez R1 G0/1 (vers R2), relancez \texttt{show ip route} et le traceroute après convergence, puis réactivez.
\question{Quel chemin de secours OSPF utilise-t-il~? Pourquoi~?}
\begin{bilan}OSPF~: famille \makebox[4.5cm]{\dotfill}~; voisins découverts par \dotfill\\Métrique \makebox[4.5cm]{\dotfill}~; aire configurée dans ce TP \makebox[4cm]{\dotfill}\end{bilan}

% ---------------------------------------------------------------- Bilan
\clearpage\section{Bilan final}
\begin{center}\begin{tabularx}{\linewidth}{lXXX}\toprule
Critère & Statique & RIPv2 & OSPF\\\midrule
Apprentissage & Manuel & Annonces de routes & État de liens\\
Choix de chemin & Route configurée & Nombre de sauts & Coût\\
Code Cisco & S & R & O\\
Adaptation aux pannes & Secours à prévoir & Convergence RIP & Convergence OSPF\\\bottomrule\end{tabularx}\end{center}
\textbf{Questions de vérification~:}
\begin{enumerate}
\item Pourquoi une route statique n'est-elle pas automatiquement mise à jour lorsqu'un nouveau réseau est créé~?\lignes
\item Un protocole dynamique peut-il restaurer la communication sans chemin physique de secours~?\lignes
\item Dans RIP, que signifie une métrique de 16~?\lignes
\item Dans OSPF, pourquoi un trajet à deux liaisons peut-il être préféré à un trajet direct~?\lignes
\item Quelle commande permet d'observer les voisins OSPF~?\lignes
\end{enumerate}
\end{document}
